\section{Further experiments: the direct macro levers}
\label{sec:experiments}

The bridge of Section~\ref{sec:bridge} covers reforms whose incidence is household-borne. This section works three experiments that instead pull the model's exogenous levers directly, chosen to exercise three different transmission channels: a pure demand-side expenditure shock, a pure price shock, and a cost-of-capital shock. All are run with \texttt{run\_reform} over 2025Q1--2027Q4 with the shock applied for twelve quarters; Table~\ref{tab:experiments} reports annual averages of the quarterly deviations, and the underlying quarterly paths are written to \texttt{figures/experiments\_quarterly.csv} by the script \texttt{figures/run\_experiments.py} that generated them.

\begin{table}[ht]
\centering
\caption{Three policy experiments through the direct levers, annual averages of quarterly deviations from the anchored baseline}
\label{tab:experiments}
\begin{tabular}{lrrrr}
\toprule
 & \multicolumn{1}{c}{$\Delta$GDP, \% } & \multicolumn{1}{c}{$\Delta$GDP} & \multicolumn{1}{c}{$\Delta$Consumption} & \multicolumn{1}{c}{$\Delta$Investment} \\
 & \multicolumn{1}{c}{of GDP} & \multicolumn{1}{c}{\pounds bn/qtr} & \multicolumn{1}{c}{\pounds bn/qtr} & \multicolumn{1}{c}{\pounds bn/qtr} \\
\midrule
\multicolumn{5}{l}{\emph{A. Government consumption $+$\pounds 5bn/yr} ($CGG$ $+$\pounds 1.25bn/qtr)}\\
\quad 2025 & $+0.18$ & $+1.250$ & $0.000$ & $0.000$ \\
\quad 2026 & $+0.19$ & $+1.250$ & $-0.000$ & $0.000$ \\
\quad 2027 & $+0.19$ & $+1.249$ & $-0.001$ & $0.000$ \\
\multicolumn{5}{l}{\emph{B. VAT-style price shock} (CPI level $+1$ per cent, $+1.4$ index points)}\\
\quad 2025 & $-0.15$ & $-1.014$ & $-1.014$ & $0.000$ \\
\quad 2026 & $-0.41$ & $-2.737$ & $-2.737$ & $0.000$ \\
\quad 2027 & $-0.46$ & $-3.070$ & $-3.070$ & $0.000$ \\
\multicolumn{5}{l}{\emph{C. Corporation tax $+$2pp} ($TCPRO$ $+0.02$, investment closure)}\\
\quad 2025 & $-0.005$ & $-0.034$ & $-0.000$ & $-0.034$ \\
\quad 2026 & $-0.04$ & $-0.295$ & $-0.004$ & $-0.290$ \\
\quad 2027 & $-0.13$ & $-0.901$ & $-0.014$ & $-0.880$ \\
\bottomrule
\end{tabular}

\medskip
\parbox{0.92\textwidth}{\footnotesize \emph{Notes:} Each experiment solves a structurally identical baseline and shocked pair (add-factors shared) and reports the delta, 2025Q1--2027Q4, shock applied for 12 quarters. Experiment C uses the investment closure with the reconstructed $\Delta \ln IBUSX$ equation~\eqref{eq:ibus} and the stabilised market-sector feedback (Section~\ref{sec:method}). Source: \texttt{figures/run\_experiments.py}; quarterly detail in \texttt{figures/experiments\_quarterly.csv}.}
\end{table}

\paragraph{A. Government consumption $+$\pounds 5bn a year: an accounting identity, not a multiplier.} The canonical demand shock: real $CGG$ rises \pounds 1.25bn per quarter. Under the demand closure the impact multiplier is one \emph{by construction}, and the label is exact rather than approximate --- $\Delta GDPM$ equals the shock to the last displayed digit. Re-running the experiment at \pounds 0.1bn and at \pounds 10bn per quarter returns an impact multiplier of 1.000000 in both cases and 0.999489 at quarter twelve in both, flat and scale-invariant: there is no behavioural content in it at all. The reported consumption response confirms this rather than qualifying it. On the \pounds 1{,}250m shock, $\Delta CONS$ peaks at $+$\pounds 0.19m, turns negative from quarter~6, reaches $-$\pounds 0.96m in 2027Q4, and cumulates to $-$\pounds 4.3m over the twelve quarters --- that is, four \emph{million} pounds against a fifteen \emph{billion} pound cumulative stimulus. An earlier version of this paragraph described that crowding-out as ``$-$\pounds 1bn cumulative by 2027Q4'', which read the final quarter's \pounds m figure as \pounds bn and overstated the response by more than two orders of magnitude; the corrected figure makes the point the wrong one obscured, which is that nothing happens.

The three channels that would generate a multiplier different from one are structurally absent, not merely small: imports have no behavioural equation in this configuration, so the leakage of equation~\eqref{eq:mgtfe} is inoperative; the income-to-consumption chain is unpopulated; and there is no monetary rule or output gap. The comparison with the OBR is therefore not the reassuring bracket this paragraph used to claim. The OBR's published multiplier for \emph{current} departmental spending is 0.6, decaying to zero over about five years \citep{obr2010multipliers}; 1.0 is its \emph{capital} multiplier. A flat, permanent, scale-invariant 1.0 is not ``between'' 0.6 and older Keynesian estimates above one --- it is the arithmetic of adding a term to an identity. Producing 0.6-with-decay would require writing equations the published listing does not contain, so this is documented rather than fixed, and the number should not be quoted as a fiscal multiplier.

\paragraph{B. A VAT-style price shock.} Because the CPI is an exogenous conditioning path (Section~\ref{sec:model}), a one per cent rise in the consumer price level ($+1.4$ index points held for twelve quarters) is a legitimate lever, and it is the closest supported analogue to the price effect of an indirect-tax rise: the OBR's costing convention for VAT changes assumes passthrough into consumer prices, which then erode real incomes. The transmission here is the deflator chain $CPI \to PCE \to RHHDI$ into the consumption ECM~\eqref{eq:cons}: with nominal income initially unchanged, real income falls one per cent, and consumption follows---$-$0.15 per cent of GDP in 2025, deepening to $-$0.46 per cent ($-$\pounds 3.1bn per quarter) by 2027 as the error-correction term drags the consumption \emph{level} toward its lower long-run solution. The entire GDP effect is consumption ($\Delta CONS = \Delta GDP$ to the pound in Table~\ref{tab:experiments}), which is the configuration speaking: investment, trade, and the labour-demand response that would spread the shock are passthrough or exogenous here. For scale, HMRC's ready reckoner puts the \emph{direct revenue} of a 1pp VAT standard-rate rise at \pounds 8.8--9.6bn a year \citep{hmrc2025reckoner}; a full VAT scoring would put that costing through the bridge as an $HHDI$-equivalent real-income loss, and the experiment here isolates the complementary price-channel arithmetic on consumer demand.

\paragraph{C. Corporation tax $+$2pp, and why the bridge refuses it.} By design, the bridge declines corporation-tax reforms: in the microsimulation, corporation tax has no household-resolution incidence---attributing it to households would smuggle in an incidence assumption (shareholders? workers? consumers?) that PolicyEngine's UK model deliberately does not make---so translating a corporation-tax costing into an $HHDI$ shock would be a category error, and the correct route is the model's own lever. Raising $TCPRO$ by two percentage points with the investment closure active works the user-cost channel $TCPRO \to TAF \to COC \to KSTAR \to KGAP \to IBUSX$: a higher tax rate raises the post-tax cost of capital, lowers desired capital, opens a capital-stock gap, and depresses business investment through the error-correction term of equation~\eqref{eq:ibus}. The response is small and slow---$-$\pounds 0.03bn per quarter in year one, building to $-$\pounds 0.9bn per quarter ($-$0.13 per cent of GDP) by 2027---consistent in sign and magnitude class with the OBR's treatment of corporation-tax measures as chiefly a receipts effect with modest medium-term investment drag.

\emph{The response converges, but slowly, and twelve quarters is well short of it.} Earlier versions of this paper reported that the deviation did not converge at all: with the closure's add-factors held as \emph{level} anchors, $|\Delta IF|$ for a sustained $+$5pp rise compounded by a factor of 1.21--1.27 every quarter, reaching \pounds 2{,}876m at quarter~12 and \pounds 43{,}387m at quarter~25 with no sign of settling. That was a bug in the anchoring, not a property of the published equation. Holding the anchors in log space instead --- the EViews convention for a \texttt{dlog} equation --- lets the published equation's own error-correction term, $-0.0418\,(\log IBUSX_{t-1} - \log KSTAR_{t-2})$, give the deviation a steady state at $\mathrm{dlog}(KSTAR) = -0.4\,\mathrm{dlog}(TAF)$. Measured on the corrected closure for the same sustained $+$5pp rise, $|\Delta IF|$ runs \pounds 0.24bn at quarter~8, \pounds 0.38bn at quarter~12 and \pounds 0.68bn at quarter~25, approaching a plateau of about \pounds 0.95bn per quarter. The error-correction root is slow ($\approx 0.958$ per quarter), so the twelve-quarter horizon every published result here uses captures only about \textbf{40 per cent} of the full effect; each run now reports its \texttt{investment\_closure\_plateau\_fraction} so a partial response is not read as the whole one. Read the sign and the direction of the corporation-tax result, and treat a twelve-quarter magnitude as a lower bound rather than the settled one. For comparison, the \emph{direct} Exchequer yield of $+$2pp on the main rate is about \pounds 7.2bn in 2026--27 rising to \pounds 8.0bn in 2028--29 on HMRC's ready reckoner \citep{hmrc2025reckoner}; the emulator's contribution is the second-round path, not the static number. The response is hard-gated in CI for sign and boundedness (Table~\ref{tab:anchored}), and Section~\ref{sec:limitations} flags the stabilised closure it depends on as a stop-gap.

\paragraph{D. The fourth lever: government investment is dead.} Section~\ref{sec:results} lists nominal government investment $CGIPS$ among the direct levers. It does not work, and the failure is silent rather than loud. Shocking $CGIPS$ by \pounds 3bn per quarter for twelve quarters leaves $\Delta IF$ at \emph{exactly} 0.0 in all twelve: $IF$ has no live equation in the published listing, so the $CGIPS \to GGIPS \to GGI \to IF$ chain never reaches the expenditure identity. What does reach GDP is deflator residue of indeterminate sign, and here it comes out \emph{negative} --- $-$0.005 to $-$0.021 per cent of GDP after a positive impact quarter --- against a published capital multiplier of 1.0 \citep{obr2010multipliers}. A user reading only the GDP column would conclude that public investment contracts the economy. The model repository now raises a warning on every $CGIPS$ shock stating that the result is not a public-investment multiplier; this paper reports no $CGIPS$ experiment because there is no response to report.

\paragraph{The scorecard for the levers.} Of the four fiscal instruments, exactly one contains behaviour. Government consumption is an accounting identity; government investment is dead and its residue is wrong-signed; corporation tax transmits but never converges. The one channel with a genuine response is the household-tax route of Section~\ref{sec:results}, whose implied multiplier of 0.09 on impact rising to 0.25 by 2027Q4 is the right magnitude class against the OBR's 0.3 income-tax assumption, though with the profile inverted: continued to twelve quarters it keeps building, passing 0.3 around quarter seven and reaching roughly 0.37, where the OBR's assumption decays. So the earlier summary of this section---``where a channel is live, the emulator reproduces the OBR architecture's sign, shape, and magnitude class; where a channel is passthrough, the emulator reports zeros and says why''---was too generous in both halves. Three of the four levers are not live in the sense that sentence implies, and two of them do \emph{not} report zeros: they report a flat 1.0 and a wrong-signed residue, both of which read as results. Labelling those at the point of use, as above and in Table~\ref{tab:anchored}, is the correction.
